\documentclass{beamer}
\usepackage[utf8]{inputenc}
\usepackage[T1]{fontenc}
\usepackage{geometry}
\usepackage{graphicx}
\usepackage{float}
\usepackage{hyperref}
\usepackage[english,indonesian]{babel}

\title{Optimasi Pemrosesan Template LaTeX Menggunakan Pendekatan Berbasis AI}
\author{Academic Support Agent}
\date{\today}

\begin{document}
\maketitle

\begin{abstract}

\end{abstract}

\chapter{PENDAHULUAN}


\chapter{TINJAUAN PUSTAKA}


% Original content (commented out):
% \documentclass[final, xcolor={table,dvipsnames},t]{beamer}

% ====================
% Packages
% ====================
\usepackage[T1]{fontenc}
\usepackage{lmodern}
\usepackage[orientation=portrait,size=a0,scale=1.0]{beamerposter} % Current dimensions A0, put in your poster dimensions
\usetheme[Blue]{UVic} % Options are DarkBlue, Blue, or White, all themed with UVic Edge styling
\usepackage{graphicx}
\usepackage{caption}
\usepackage{booktabs}
\usepackage{tikz}
\usepackage{pgfplots}
\pgfplotsset{compat=1.17}
\newcommand{\blu}{\color{blue}}


\usepackage{blindtext}


%%%%%%%%%%%%%%%%%%%%%%%%%%%%%%%%%%%%%%%%%%%%%%%%%%%%%%%%%%%%%%%%%%%%%%%%%%%%%%
% Column environment setup
%%%%%%%%%%%%%%%%%%%%%%%%%%%%%%%%%%%%%%%%%%%%%%%%%%%%%%%%%%%%%%%%%%%%%%%%%%%%%%
% If you have N columns, choose \sepwidth and \colwidth such that
% (N+1)*\sepwidth + N*\colwidth = \paperwidth
% Follow structure to create difference column environments. 
\newlength{\sepwidthA} % Separation distance between columns type A
\newlength{\colwidthA} % column width type A
\setlength{\sepwidthA}{0.25\paperwidth}
\setlength{\colwidthA}{0.5\paperwidth}

\newcommand{\separatorcolumnA}{\begin{column}{\sepwidthA}\end{column}}

% Second column environment. 
\newlength{\sepwidthB}
\newlength{\colwidthB}
\setlength{\sepwidthB}{0.0666\paperwidth}
\setlength{\colwidthB}{0.4\paperwidth}

\newcommand{\separatorcolumnB}{\begin{column}{\sepwidthB}\end{column}}

% You can also use these column commands to create columns inside columns and for creating new column formatting. 
% You can also have non even columns by creating more column environments or specifying the width when beginning a column environment. 
%%%%%%%%%%%%%%%%%%%%%%%%%%%%%%%%%%%%%%%%%%%%%%%%%%%%%%%%%%%%%%%%%%%%%%%%%%%%%%
% Title
%%%%%%%%%%%%%%%%%%%%%%%%%%%%%%%%%%%%%%%%%%%%%%%%%%%%%%%%%%%%%%%%%%%%%%%%%%%%%%

\title{\VeryHuge{Centered Title}}
\author{\textbf{Presenting Author}, Author Two and Author Three}
\institute[shortinst]{Author Institution, Place of Institution}

%%%%%%%%%%%%%%%%%%%%%%%%%%%%%%%%%%%%%%%%%%%%%%%%%%%%%%%%%%%%%%%%%%%%%%%%%%%%%%
% Poster footer
%%%%%%%%%%%%%%%%%%%%%%%%%%%%%%%%%%%%%%%%%%%%%%%%%%%%%%%%%%%%%%%%%%%%%%%%%%%%%%

\footercontent{
\href{mailto:Presenting.Author@uvic.ca}{Presenting.Author@uvic.ca} % this is a clickable link
  \hfill
  Conference Name \hfill
  {Poster Number: \#} }
% (can be left out to remove footer)    

\begin{document}
%%%%%%%%%%%%%%%%%%%%%%%%%%%%%%%%%%%%%%%%%%%%%%%%%%%%%%%%%%%%%%%%%%%%%%%%%%%%%%
% Logo placements (optional)
%%%%%%%%%%%%%%%%%%%%%%%%%%%%%%%%%%%%%%%%%%%%%%%%%%%%%%%%%%%%%%%%%%%%%%%%%%%%%%

% \addtobeamertemplate{headline}{}
% {
%     \begin{tikzpicture}[remember picture,overlay,line width=\arrayrulewidth]
%     %   Logo in top right below header
%       \node [anchor=north east, inner sep=3cm] at ([xshift=1.0cm,yshift=-10.0cm]current page.north east)
%       {\includegraphics[height=6.0cm]{logos/mylogo.eps}}; 
%         % Logo in top left below header
%       \node [anchor=north west, inner sep=3cm] at ([xshift=-1.0cm,yshift=-10.0cm]current page.north west)
%       {\includegraphics[height=7.0cm]{logos/mylogo.eps}}; 
%     \end{tikzpicture}
% }

% ====================
% Body
% ====================

\begin{frame}[t]
%%%%%%%%%%%%%%%%%%%%%%%%%%%%%%%%%%%%%%%%%%%
%Section 1
%%%%%%%%%%%%%%%%%%%%%%%%%%%%%%%%%%%%%%%%%%%
\begin{columns}[t]
    \separatorcolumnA
    \begin{column}{\colwidthA}

        \begin{block}{Introduce graphical abstract, pose research question}
            \blindtext
        \end{block}
    
    \end{column}

    \separatorcolumnA
\end{columns}

%%%%%%%%%%%%%%%%%%%%%%%%%%%%%%%%%%%%%%%%%%%
%Section 2 
%%%%%%%%%%%%%%%%%%%%%%%%%%%%%%%%%%%%%%%%%%%
\begin{columns}
%%%%%%%%%%%%%%%%%%%%%%%%%%%%%%%%%%%%%%%%%%%
%Section 2 column 1
%%%%%%%%%%%%%%%%%%%%%%%%%%%%%%%%%%%%%%%%%%%
\separatorcolumnB

    \begin{column}[T]{\colwidthB}

        \begin{block}{Block with a diagram/figure}

        Here is a diagram of a star drawn in tikz.

        \begin{figure}
            \centering
                \begin{tikzpicture}[scale=8]
                \draw[step=0.25cm,color=gray] (-1,-1) grid (1,1);
                \draw (1,0) -- (0.2,0.2) -- (0,1) -- (-0.2,0.2) -- (-1,0)
                -- (-0.2,-0.2) -- (0,-1) -- (0.2,-0.2) -- cycle;
                \end{tikzpicture}
            \caption{A figure caption. If needed for data or reference}
        \end{figure}
    \end{block}
    
    \begin{block}{Block with nested columns}
        \begin{column}{0.3\colwidthB}
            \begin{itemize}
                \item Keep text information clear and concise
                \item Include some math as well: $ F = ma $
            \end{itemize}
        \end{column}
        \begin{column}{0.7\colwidthB}
            Some possible justification for having sub-columns and why you shouldn't be banished to the nether.
            
            Also some conversation on the math out left because despite this formatting travesty it is an interesting point. Does $F=ma$?
        \end{column}
    \end{block}

    \begin{alertblock}{Highlighted Block}
        Use highlighting to bring attention to a key section or result, can also be used to reduce contrast with images on white background. 
    \end{alertblock}
        
\end{column}
\separatorcolumnB
%%%%%%%%%%%%%%%%%%%%%%%%%%%%%%%%%%%%%%%%%%%
%Section 2 column 2
%%%%%%%%%%%%%%%%%%%%%%%%%%%%%%%%%%%%%%%%%%%
\begin{column}[T]{\colwidthB}

    \begin{block}{Block With a Table}
    A table in a block for torturing your audience. 
        \begin{table}
            \centering
                \begin{tabular}{l c c c}
                \tableheadrow
                \tableheadcol{ Team } & \tableheadcol{ Win } & \tableheadcol{ Loss } & \tableheadcol{ Tie } \\
                Canucks & 10 & 0 & 0 \\
                Leafs & 0 & 9 & 1 \\
                Canadians & 7 & 3 & 0 \\
                \bottomrule
            \end{tabular}
            \caption{A table caption.}
        \end{table}
    \end{block}

  

    \begin{block}{QR code for linking online resources}
        \begin{figure}
            \centering
            \includegraphics[width=0.7\textwidth]{figures/dololphin.jpg}
        \end{figure}
\begin{column}[T]{0.6\colwidthB}
      This QR code goes to my poster online for your convenience. Also use it to check out our publications and other research. 
\end{column}
\begin{column}[T]{0.4\colwidthB}
    \begin{figure}
      \centering
      \includegraphics[width=0.5\textwidth]{figures/LearnMore.eps}
      \captionsetup{labelformat=empty}
        \caption{\href{https://www.youtube.com/watch?v=dQw4w9WgXcQ}{\textbf{\blu{Clickable link to URL}}}}
     \end{figure}
\end{column}
\end{block}

%%%%%%%%%%%%%%%%%%%%%%%%%%%%%%%%%%%%%%%%%%%%%%%%%%%%%%%%%%%%%%%%%%%%%%%%%%%%%%
% References
%%%%%%%%%%%%%%%%%%%%%%%%%%%%%%%%%%%%%%%%%%%%%%%%%%%%%%%%%%%%%%%%%%%%%%%%%%%%%%
\begin{block}{References}

\nocite{*}
\bibliography{poster}% Produces the bibliography via BibTeX.
\bibliographystyle{plain}
\end{block} 

\end{column}
\separatorcolumnB
\end{columns}
\end{frame}
\end{document}


\end{document}